\section{Restriction Categories and Categories with Support}

For the theory of data, 
the notion of \textit{range category} that we are interested in 
are define in terms of restriction categories.

\subsection{Restriction Categories}
These were introduced by 
Cockett \& Lack (\cite{COCKETT2002}) and 
constitute an algebra that nicely represents a system of partially defined terms or maps. Applied to the theory of data, they are relevant because data fields often only optionally have values and can be considered therefore to be partial, rather than total, functional relationships. 

Cockett \& Lack define restriction categories along the following lines:
\begin{definition}
A \textit{restriction category} is a category \catcw with an additional operator on morphisms (the restriction operator $\rst{\ }$), so that from a morphism
$f: a \morph b$ in \catcw there is a 
$$\rst{f}: a \morph a$$,
satisfying

R.1 For $f:a \morph b$ in \catcw $$\rst{f} \circ f =f$$.

R.2. If \fgsourcediag in \catcw then
$$\rst{g} \circ \rst{f}=\rst{f} \circ \rst{g}.$$

R.3. If \fgsourcediag in \catcw then
$$\rst{\rst{f} \circ g} = \rst{f} \circ \rst{g}$$.

R.4. If $\sequentialdiag{a}{b}{c}{f}{g}$ in \catcw then
$$f \circ \rst{g} = \rst{f \circ g} \circ f$$.
\end{definition}

We will be needing:

\begin{lemma}
\llabel{restrictioncatlemma}
In any restriction category \catc
\begin{enumerate} [(i)]
\item For $f:a \morph b$ in \catcw,
$$\rst{\rst{f}}=\rst{f}$$
\item w.R4.\footnote{This condition described as wR.4 in Cockett et al.} 
If $\sequentialdiag{a}{b}{c}{f}{g}$ in \catcw 
 then
$$\rst{f \circ g} = \rst{f \circ \rst{g}}.$$

\item \llabel{CHECKLemma}
If \catcw is a restriction category \catcw and if
$\sequentialdiag{a}{b}{c}{f}{g}$ in \catcw then 
$f \circ \rst{g} = f$ iff $\rst{f \circ g} = \rst{f}.$
\item \llabel{restrictioncatlemma:totalk} \commentary{UNUSED}
If $\sequentialdiag{a}{b}{c}{k}{j}$ in \catcw 
and $k$ is total i.e. $\rst{k}=id_a$
and if $k \circ \rst{j}$ is total i.e. $\rst{k \circ \rst{j}}=id_a$
then
$$k \circ \rst{j} = k.$$
\end{enumerate}
\end{lemma}
\begin{proof}
\begin{enumerate} [(i)]
\item Use R3 with $g$ being $id_b$.
\item 
\begin{align*}
\rst{f \circ \rst{g}}
&= \rst{\rst{f \circ g} \circ f}
    && \text{by (R4),} \\
&= \rst{f \circ g} \circ \rst{f}
    && \text{by (R3),} \\
&= \rst{f} \circ \rst{f \circ g}
    && \text{by (R2),} \\
  &= \rst{\rst{f} \circ f \circ g}
    && \text{by (R3),} \\  
&= \rst{f \circ g}
    && \text{by (R1).}
\end{align*}
\item \commentary{CHECKLemma}
LEFT TO RIGHT
Assume $f \circ \rst{g} = f$
then
\begin{align*}
\rst{f \circ g} &= \rst{f \circ \rst{g}} && \mbox{by wR.4,}\\
                     &= \rst{f} && \mbox{from the assumption.}
\end{align*}
RIGHT TO LEFT
Assume $\rst{f \circ g} = \rst{f}$
then
\begin{align*}
f \circ \rst{g} &= \rst{f \circ g} \circ f 
                                  && \mbox{R.4,} \\
                &= \rst{f} \circ f 
                                  && \mbox{from the assumption,} \\
                &= f              && \mbox{by R.1.}
\end{align*}
\item 
\begin{align*}
k \circ \rst{j} &= k \circ \rst{\rst{j}}              && \mbox{from (i),}\\
                &= \rst{k \circ \rst{j}} \circ k && \mbox{by R.4,}   \\
                &= k                                  && \mbox{from assumption $k \circ \rst{j}$ is total.}
\end{align*}

\end{enumerate}
\end{proof}

\subsection{Categories with Support}
We are  interested in restriction categories because they are the basis of  range categories. 
But to get to  where we need to be we need to work with 
\textit{categories with support}. I use the term as it used
in Cockett \cite{COCKETT2012}. In summary a category with support is very much like a restriction category. It has an operator  \rstItself for which R.1, R.2 and R.3 hold but for which R.4 does not necessarily hold. Instead the following
weaker condition holds

wR.4  If $\sequentialdiag{a}{b}{c}{f}{g}$ in \catcw then
\[
\rst{f \circ g} = \rst{f \circ \rst{g}}.
\]
Such categories Cockett et al refer to as \textit{categories with support}. 

Categories with support may satisy an addition condition which is crucial in 
what follows and which we refer to as the additional weak R.4 rule:

awR.4  If $\sequentialdiag{a}{b}{c}{f}{g}$ in \catcw then
\[
\rst{f \circ g} = \rst{f} \mbox{ \quad implies \quad} f \circ \rst{g}=f.
\]

\iffalse UNUSED (NO LONGER USED)
***********************************
\begin{lemma}
\llabel{supportLemma}
In any category with support \catc, 
\begin{enumerate}[(i)]

\item
if $\sequentialdiag{a}{b}{c}{f}{g}$ in \catcw 
 then
 \[
 \rst{f} \circ \rst{f \circ g} = \rst{f \circ g}
 \]
 \item \llabel{RHomFunctorSubLemma}
if  $\triplesequentialdiag{a}{x_1}{x_2}{x_3}{k}{j_1}{j_2}$ 
in a \catcw  then 
if $\rst{k \circ j_1 \circ j_2} = \rst{k}$ 
then $\rst{k \circ j_1} = \rst{k}$.
\end{enumerate}
\end{lemma}
\begin{proof}
\begin{enumerate}[(i)]
\item
\begin{align*}
\rst{f} \circ \rst{f \circ g}  
           &= \rst{f} \circ \rst{f \circ \rst{g}}
                               && \mbox{by wR.4} \\
           &= \rst{\rst{f} \circ f \circ \rst{g}}
                               && \mbox{by R.3}  \\
           &= \rst{f \circ \rst{g}}
                               && \mbox{by R.1}\\
           &= \rst{f \circ g}
                               && \mbox{ by wR.4, as required}
\end{align*}
\item

Suppose $\rst{k \circ j_1 \circ j_2} = \rst{k}$  then we have
\begin{align*}
\rst{k \circ j_1} &= \rst{\rst{k} \circ k \circ j_1}
                           &&\mbox{by R.1,} \\
                       &= \rst{\rst{k \circ j_1 \circ j_2} \circ k \circ j_1}
                           &&\mbox{from the initial assumption,} \\
                       &= \rst{k \circ j_1 \circ j_2} \circ \rst{k \circ j_1}
                           &&\mbox{by R.3,} \\
                       &= \rst{k \circ j_1 \circ j_2}
                           &&\mbox{by (i),} \\
                       &= \rst{k}
                           &&\mbox{from the initial assumption, as required.} 
\end{align*}
\end{enumerate}
\end{proof}
***************************
\fi
\begin{lemma}
\llabel{newSupportLemma}
If \fgsourcediag in a category with support \catcw then
\[
\rst{\rst{g} \circ f} \circ g = \rst{f} \circ g
\]
\end{lemma}
\begin{proof}
\begin{align*}
\rst{\rst{g} \circ f} \circ g
    &= \rst{g} \circ \rst{f} \circ g && \mbox{by R.3,}\\
    &= \rst{f} \circ \rst{g} \circ g  && \mbox{by R.2,}\\
    &= \rst{f} \circ g               && \mbox{by R.1.}
\end{align*}
\end{proof}

\subsection{Partial Ordering of $Hom(A,B)$}

In a category with support, and therefore in a restriction category, we can define a partial ordering on each hom set
Hom(a,b) by defining :
$$f \leq g \mbox{ iff } f = \rst{f} \circ g.$$

The relation $\leq$, so defined, is antisymmetric because if
\begin{equation}
\label{fleqg}
f = \rst{f} \circ g
\end{equation}
and
\begin{equation}
\label{gleqf}
g = \rst{g} \circ f
\end{equation}
then
\begin{align*}
f &= \rst{f} \circ g &&\mbox {this is (\ref{fleqg}),}        \\
  &= \rst{g} \circ \rst{f} \circ g &&\mbox{by R.1 and R.2,}\\
  &= \rst{g} \circ f               &&\mbox{using (\ref{fleqg}) again,} \\
  &= g                             &&\mbox{(by (\ref{gleqf}).} 
\end{align*}

When $f$ and $g$ are partial functions and
considered as morphisms in the category $\Par$ of sets and partial functions,  
then  $f \leq g$ iff  whenever $f$ is defined at an $x$,  then $g$ too is defined at x and $f(x) = g(x)$.

\begin{lemma}
\llabel{orderingsCompose}
In a restriction category \catcw, \commentary{category with support ??}
if 
$
\paralleldiag{a}{b}{f}{f'}
$
and
$
\paralleldiag{b}{c}{g}{g'}
$
in \catcw and if $f' \leq f$ and $g' \leq g$ 
then $f' \circ g' \leq f \circ g$.
\end{lemma}

\begin{lemma}
\llabel{pointwiseorderingsCompose}
In a restriction category \catcw, \commentary{category with support ??}
if 
$
\paralleldiag{a}{b}{f}{f'}
$
\nudgedown{0.6cm}
and
$
\paralleldiag{b}{c}{g}{g'}
$
and $h: a \morph c$ in \catcw 
and if $f' \leq f$ and $g' \leq g$ 
and if $h \leq f \circ g'$ and $h \leq f' \circ g$ 
then $h \leq f' \circ g'$.
\end{lemma}
\begin{proof}
Proof is handwritten on tablet.\highlight{CHECK.}
\end{proof}


\begin{lemma}
\llabel{leastFactorLemma}
\commentary{19 May 2026: \workt{Need} an n-ary version of this}
\commentary{19 May 2026: applies to referential inclusion dependencies}

\[
\begin{array}{c p{0.25cm} c p {0.25cm} c}
       &&         && \Ob{c}{1}  \\[0.4cm]
\bOb{a}&& \bOb{b} &&            \\[0.4cm]
       &&         && \Ob{c}{2} 
\end{array}
\begin{arrows}
\ncarr{a}{b}\alabel{f}
\ncarr{b}{c1}\alabel{q_1}
\ncarr{b}{c2}\blabel{q_2}
\end{arrows}       
\]
in a restriction category \catcw then
if $\rst{f \circ q_1}=\rst{f \circ q_2}$
then 
\begin{equation}
\label{leastFactorOne}
f \circ \rst{q_1} \circ \rst{q_2} \circ q_1=f \circ q_1
\end{equation}
(and by symmetry )
\begin{equation}
\label{leastFactorTwo}
f \circ \rst{q_1} \circ \rst{q_2} \circ q_2=f \circ q_2
\end{equation}
\end{lemma}
\begin{proof}

\begin{align*}
f \circ \rst{q_1} \circ \rst{q_2} \circ q_1 
    &= f \circ \rst{q_2} \circ q_1 
       && \mbox{using R.2 and R.1}, \\
    &= \rst{f \circ q_2} \circ f \circ q_1
       && \mbox{by R.4,} \\
    &= \rst{f \circ q_1} \circ f \circ q_1
       && \mbox{since 
      $\rst{f \circ q_1} = \rst{f \circ q_2}$ } \\
    &= f \circ q_1
       && \mbox{by R.1, as required.}
\end{align*} 
\end{proof}

